% Source: https://artofproblemsolving.com/community/c3h1567541_fractions
% Stable identity: aops-topic-1567541-q01
First $\dfrac{68}{119}=\dfrac{4\cdot17}{7\cdot17}=\dfrac{4}{7}$. Both $a$ and $b$ are nonzero digits (they lead two-digit numbers), and
\[
\dfrac{10a+b}{10b+a}=\dfrac{4}{7}\iff 70a+7b=40b+4a\iff 66a=33b\iff b=2a .
\]
So $a\in\{1,2,3,4\}$ (as $b=2a\le9$), giving $\dfrac{12}{21},\dfrac{24}{42},\dfrac{36}{63},\dfrac{48}{84}$, each equal to $\dfrac{4}{7}$. There are $\boxed{4}$ such fractions.
